\documentclass[11pt]{article}
\usepackage[T1]{fontenc}
\usepackage[margin=1in]{geometry}
\usepackage{amsmath,amssymb}
\usepackage[bookmarks=true,bookmarksnumbered=true,colorlinks=true,linkcolor=blue,urlcolor=blue]{hyperref}

\title{Energy Methods in Classical Mechanics}
\author{A. Reyes}
\date{}

\begin{document}
\maketitle
\tableofcontents

\section{Introduction}
\label{sec:intro}

This note collects the two variational principles that organize classical
mechanics: the Lagrangian principle of Section~\ref{sec:lagrange} and the
Hamiltonian formulation of Section~\ref{sec:hamilton} on
page~\pageref{sec:hamilton}. Both describe the same trajectories, but they
generalize in different directions, and the choice between them matters as
soon as constraints or symmetries enter the picture.

The running example is the planar pendulum of mass $m$ and length $\ell$,
with generalized coordinate $\theta$ measured from the downward vertical.
Its kinetic and potential energies are
\begin{equation}
  \label{eq:energies}
  T = \tfrac12 m \ell^{2} \dot{\theta}^{2}, \qquad
  V = -m g \ell \cos\theta,
\end{equation}
and every abstract statement below can be tested against this system.

\section{Lagrangian mechanics: the \texorpdfstring{$L = T - V$}{L = T - V} principle}
\label{sec:lagrange}

For a system with generalized coordinates $q = (q_1, \dots, q_n)$ define the
Lagrangian $L(q, \dot{q}, t) = T - V$. Hamilton's principle states that the
physical path renders the action $S = \int L \, dt$ stationary, which yields
the Euler--Lagrange equations
\begin{equation}
  \label{eq:euler-lagrange}
  \frac{d}{dt}\frac{\partial L}{\partial \dot{q}_i}
    - \frac{\partial L}{\partial q_i} = 0,
    \qquad i = 1, \dots, n.
\end{equation}
Applied to the pendulum of \eqref{eq:energies}, equation
\eqref{eq:euler-lagrange} gives $\ddot{\theta} + (g/\ell)\sin\theta = 0$,
the familiar nonlinear equation of motion.

\subsection{Euler--Lagrange equations with \texorpdfstring{$\dot{q}$}{q-dot} dependence}
\label{sec:qdot}

The velocity dependence of $L$ is what makes the theory dynamical rather than
merely geometrical: the Hessian matrix $\partial^2 L / \partial \dot{q}_i
\partial \dot{q}_j$ must be invertible for the accelerations to be determined
uniquely by positions and velocities. When this Hessian is singular, the
system is constrained and must be handled by the Dirac--Bergmann algorithm,
which properly belongs to the Hamiltonian side described in
Section~\ref{sec:hamilton}.

\subsubsection{Dissipation via the \texorpdfstring{$\mathcal{F}$}{F} function}
\label{sec:dissipation}

Weak friction is accommodated by Rayleigh's dissipation function
$\mathcal{F} = \tfrac12 b \ell^{2} \dot{\theta}^{2}$, which adds a term
$-\partial \mathcal{F} / \partial \dot{q}_i$ to the right-hand side of
\eqref{eq:euler-lagrange}. The damped pendulum then satisfies
$\ddot{\theta} + \gamma \dot{\theta} + (g/\ell)\sin\theta = 0$ with
$\gamma = b/m$, and its total energy decays monotonically.

\subsection{Cyclic coordinates and conservation}
\label{sec:cyclic}

A coordinate $q_k$ is cyclic when $\partial L / \partial q_k = 0$; its
conjugate momentum $p_k = \partial L / \partial \dot{q}_k$ is then conserved.
This is the Lagrangian shadow of Noether's theorem, whose full statement is
cleanest in the Hamiltonian variables introduced next.

\section{Hamiltonian formulation with \texorpdfstring{$H(p, q)$}{H(p,q)}}
\label{sec:hamilton}

The Legendre transform $p_i = \partial L / \partial \dot{q}_i$ trades
velocities for momenta and produces the Hamiltonian $H(p, q) = \sum_i p_i
\dot{q}_i - L$. The equations of motion become the first-order system
\begin{equation}
  \label{eq:hamilton-eqs}
  \dot{q}_i = \frac{\partial H}{\partial p_i}, \qquad
  \dot{p}_i = -\frac{\partial H}{\partial q_i}.
\end{equation}
For the pendulum, $p_\theta = m\ell^2\dot\theta$ and
$H = p_\theta^2/(2m\ell^2) - mg\ell\cos\theta$, and \eqref{eq:hamilton-eqs}
reproduces the same trajectories as \eqref{eq:euler-lagrange} while making
energy conservation, $dH/dt = 0$, manifest.

\subsection{Poisson brackets with \texorpdfstring{$\{q, p\} = 1$}{{q,p} = 1}}
\label{sec:poisson}

Observables evolve by $\dot{f} = \{f, H\}$, where the Poisson bracket of the
canonical pair satisfies $\{q_i, p_j\} = \delta_{ij}$. Quantization promotes
this bracket to the commutator, which is why the Hamiltonian picture of this
section, rather than the Lagrangian picture of Section~\ref{sec:lagrange},
is the gateway to quantum mechanics.

\end{document}
